% -*- coding: utf-8 -*-
% XeLaTeX 编译
\documentclass[a3paper, landscape, 11pt, twocolumn]{article}

% 引入宏包
\usepackage{fontspec}
\usepackage[UTF8]{ctex}
\usepackage{geometry}
\usepackage{listings}
\usepackage{xcolor}
\usepackage{enumitem}
\usepackage{tabularx}
\usepackage{makecell}
\usepackage{fancyhdr}
\usepackage{amssymb}
\usepackage{lastpage}

% 设置页面边距 (A3横向双栏)
\geometry{left=2cm, right=2cm, top=2.5cm, bottom=2.5cm, columnsep=1.8cm}

% 颜色定义
\definecolor{codebg}{RGB}{245,245,245}
\definecolor{keywordcolor}{RGB}{0,0,150}
\definecolor{commentcolor}{RGB}{0,128,0}
\definecolor{stringcolor}{RGB}{163,21,21}

% Java代码块样式设置
\lstdefinestyle{JavaExam}{
    language=Java,
    basicstyle=\ttfamily\small,
    backgroundcolor=\color{codebg},
    keywordstyle=\color{keywordcolor}\bfseries,
    commentstyle=\color{commentcolor},
    stringstyle=\color{stringcolor},
    showstringspaces=false,
    breaklines=true,
    frame=single,
    rulecolor=\color{black!50},
    tabsize=4,
    xleftmargin=1em,
    xrightmargin=1em,
    aboveskip=1em,
    belowskip=1em
}

% 列表样式
\newlist{options}{enumerate}{1}
\setlist[options]{label=\Alph*., nosep, leftmargin=*, itemsep=0.5em}

% 选择题选项命令
\newcommand{\chooseoption}{\begin{options}\item}

% 页眉页脚设置
\pagestyle{fancy}
\fancyhf{}
\fancyhead[C]{\Large\bfseries 《Java程序设计》期末考试试卷}
\fancyfoot[C]{第 \thepage\ 页 共 \pageref{LastPage}\ 页}
\fancyfoot[R]{得分：\underline{\hspace{2cm}}}
\fancyfoot[L]{题号：\underline{\hspace{1cm}}}
\fancyheadoffset{0cm}
\fancyfootoffset{0cm}
\fancyfootoffset{0cm}
\fancyfootoffset{0cm}
\fancyfootoffset{0cm}
\fancyfootoffset{0cm}
\renewcommand{\headrulewidth}{0.4pt}

\begin{document}

% 试卷装订线与密封线模拟 (绝对定位，此处用简单文本示意)
\twocolumn[{
			\begin{center}
				\vspace*{-1cm}
				{\huge\bfseries 2025—2026学年第一学期《Java程序设计》期末试卷}\\[1em]
				{\large 适用专业：计算机科学与技术、软件工程 \hspace{2em} 闭卷 \hspace{2em} 满分：150分 \hspace{2em} 考试时间：120分钟}\\[1em]
				\renewcommand{\arraystretch}{1.5}
				\begin{tabular}{|c|c|c|c|c|c|c|}
					\hline
					题号  & 一 (40分)        & 二 (20分)        & 三 (20分)        & 四 (30分)        & 五 (40分)        & 总分 (150分)      \\
					\hline
					得分  & \hspace{1.5cm} & \hspace{1.5cm} & \hspace{1.5cm} & \hspace{1.5cm} & \hspace{1.5cm} & \hspace{1.5cm} \\
					\hline
					阅卷人 &                &                &                &                &                &                \\
					\hline
				\end{tabular}\\[1.5em]
				{\large 姓名：\underline{\hspace{3cm}} \hspace{2em} 学号：\underline{\hspace{4cm}} \hspace{2em} 班级：\underline{\hspace{3cm}}}
				\vspace{1em}
				\hrule
				\vspace{1.5em}
			\end{center}
		}]

\section*{一、单项选择题（本大题共 20 小题，每小题 2 分，共 40 分）}
\textit{在每小题列出的四个备选项中只有一个是符合题目要求的，请将其代码填写在题后的括号内。错选、多选或未选均无分。}

\begin{enumerate}[label=\arabic*., itemsep=1em]
	\item Java程序的执行过程是由解释器执行编译后生成的字节码文件，该字节码文件的扩展名是（ \quad ）。
	      \begin{options}
		      \item \texttt{.java}
		      \item \texttt{.class}
		      \item \texttt{.exe}
		      \item \texttt{.jar}
	      \end{options}

	\item 在Java中，下列哪个标识符是不合法的？（ \quad ）
	      \begin{options}
		      \item \texttt{\_studentName}
		      \item \texttt{\$salary}
		      \item \texttt{2ndNumber}
		      \item \texttt{MAX\_VALUE}
	      \end{options}

	\item Java基本数据类型中，\texttt{char} 类型所占用的内存空间是（ \quad ）个字节。
	      \begin{options}
		      \item 1
		      \item 2
		      \item 4
		      \item 8
	      \end{options}

	\item 假设 \texttt{a} 为整型变量，初值为 5，则执行表达式 \texttt{a += a * 2} 后，\texttt{a} 的值为（ \quad ）。
	      \begin{options}
		      \item 10
		      \item 15
		      \item 20
		      \item 25
	      \end{options}

	\item 下列哪种数据类型不能作为 \texttt{switch} 语句的条件测试变量？（ \quad ）
	      \begin{options}
		      \item \texttt{byte}
		      \item \texttt{long}
		      \item \texttt{char}
		      \item \texttt{String}
	      \end{options}

	\item 在定义二维数组时，下列哪种写法是错误的？（ \quad ）
	      \begin{options}
		      \item \texttt{int[][] a = new int[3][4];}
		      \item \texttt{int a[][] = new int[3][];}
		      \item \texttt{int[][] a = new int[][4];}
		      \item \texttt{int[] a[] = \{\{1,2\}, \{3,4\}\};}
	      \end{options}

	\item 面向对象程序设计的三个主要特征不包括（ \quad ）。
	      \begin{options}
		      \item 封装性
		      \item 继承性
		      \item 健壮性
		      \item 多态性
	      \end{options}

	\item 关于Java类中的构造方法，下列说法正确的是（ \quad ）。
	      \begin{options}
		      \item 构造方法的方法名可以和类名不同
		      \item 构造方法必须声明返回值类型，即使是 \texttt{void}
		      \item 一个类只能有一个构造方法
		      \item 构造方法在通过 \texttt{new} 关键字实例化对象时自动调用
	      \end{options}

	\item 若要在子类的构造方法中调用父类的有参构造方法，必须使用（ \quad ）关键字，并且必须放在子类构造方法的第一行。
	      \begin{options}
		      \item \texttt{this}
		      \item \texttt{super}
		      \item \texttt{final}
		      \item \texttt{static}
	      \end{options}

	\item 下列关于 \texttt{final} 关键字的描述，错误的是（ \quad ）。
	      \begin{options}
		      \item 被 \texttt{final} 修饰的类不能被继承
		      \item 被 \texttt{final} 修饰的方法不能被重写
		      \item 被 \texttt{final} 修饰的变量即为常量，一旦赋值后不能修改
		      \item 抽象类可以被 \texttt{final} 修饰
	      \end{options}

	\item 在Java中，所有类的顶级父类是（ \quad ）。
	      \begin{options}
		      \item \texttt{java.lang.Class}
		      \item \texttt{java.lang.Object}
		      \item \texttt{java.util.Scanner}
		      \item \texttt{java.lang.System}
	      \end{options}

	\item 接口（Interface）中的成员变量默认的修饰符是（ \quad ）。
	      \begin{options}
		      \item \texttt{private static final}
		      \item \texttt{public static final}
		      \item \texttt{protected static}
		      \item 默认无修饰符
	      \end{options}

	\item 当尝试访问数组的非法索引时，Java运行时会抛出下列哪个异常？（ \quad ）
	      \begin{options}
		      \item \texttt{NullPointerException}
		      \item \texttt{ArithmeticException}
		      \item \texttt{ArrayIndexOutOfBoundsException}
		      \item \texttt{NumberFormatException}
	      \end{options}

	\item 下列哪个类主要用于处理频繁修改的字符串序列，且是线程安全的？（ \quad ）
	      \begin{options}
		      \item \texttt{String}
		      \item \texttt{StringBuilder}
		      \item \texttt{StringBuffer}
		      \item \texttt{Character}
	      \end{options}

	\item Java集合框架中，如果希望元素能够自动保持排序（如按自然顺序），应选择下列哪个实现类？（ \quad ）
	      \begin{options}
		      \item \texttt{HashSet}
		      \item \texttt{ArrayList}
		      \item \texttt{LinkedList}
		      \item \texttt{TreeSet}
	      \end{options}

	\item 在Java I/O 流中，以字符为单位进行读写操作的基础抽象类是（ \quad ）。
	      \begin{options}
		      \item \texttt{InputStream} 和 \texttt{OutputStream}
		      \item \texttt{Reader} 和 \texttt{Writer}
		      \item \texttt{FileInputStream} 和 \texttt{FileOutputStream}
		      \item \texttt{DataInputStream} 和 \texttt{DataOutputStream}
	      \end{options}

	\item 下列创建和启动多线程的语句中，正确的是（ \quad ）。
	      \begin{options}
		      \item \texttt{new Thread(new RunnableImpl()).run();}
		      \item \texttt{new Thread(new RunnableImpl()).start();}
		      \item \texttt{new RunnableImpl().start();}
		      \item \texttt{Thread.start(new RunnableImpl());}
	      \end{options}

	\item 要使某个类具备被序列化的能力，该类必须实现（ \quad ）接口。
	      \begin{options}
		      \item \texttt{Cloneable}
		      \item \texttt{Runnable}
		      \item \texttt{Serializable}
		      \item \texttt{Comparable}
	      \end{options}

	\item GUI编程中，\texttt{JFrame} 的默认布局管理器是（ \quad ）。
	      \begin{options}
		      \item \texttt{FlowLayout}
		      \item \texttt{GridLayout}
		      \item \texttt{BorderLayout}
		      \item \texttt{CardLayout}
	      \end{options}

	\item JDBC操作数据库时，用于执行预编译SQL语句的对象是（ \quad ）。
	      \begin{options}
		      \item \texttt{Statement}
		      \item \texttt{PreparedStatement}
		      \item \texttt{CallableStatement}
		      \item \texttt{ResultSet}
	      \end{options}

\end{enumerate}

\section*{二、判断题（本大题共 10 小题，每小题 2 分，共 20 分）}
\textit{请判断下列说法的正误。正确的打“$\checkmark$”，错误的打“$\times$”。}

\begin{enumerate}[label=\arabic*., itemsep=0.8em]
	\item Java中的数组一旦创建，其长度就固定不可改变。 （ \quad ）
	\item \texttt{==} 运算符用于比较两个对象的内存地址是否相同，而 \texttt{equals()} 方法在未被重写时比较的是对象的属性值。 （ \quad ）
	\item 在Java中，一个子类只能继承一个父类，但一个类可以实现多个接口。 （ \quad ）
	\item 父类的构造方法可以被子类继承和直接使用。 （ \quad ）
	\item 带有 \texttt{try-catch-finally} 的异常处理结构中，无论是否发生异常，\texttt{finally} 块中的代码通常都会被执行。 （ \quad ）
	\item \texttt{String} 类型的对象是不可变的（Immutable），一旦创建就不能修改其内容。 （ \quad ）
	\item 在多重 \texttt{catch} 块中，必须将捕获 \texttt{Exception} 类的块放在捕获 \texttt{IOException} 类的块之前。 （ \quad ）
	\item \texttt{ArrayList} 内部是基于动态数组实现的，因此它的随机访问速度比 \texttt{LinkedList} 更快。 （ \quad ）
	\item 包含抽象方法的类必须声明为抽象类，但抽象类也可以不包含任何抽象方法。 （ \quad ）
	\item 方法内部声明的局部变量如果没有进行显式初始化，系统会自动为其赋予默认值。 （ \quad ）
\end{enumerate}


\section*{三、填空题（本大题共 5 小题，每空 2 分，共 20 分）}

\begin{enumerate}[label=\arabic*., itemsep=1em]
	\item Java的跨平台特性主要依赖于Java虚拟机（JVM）。Java源代码被编译后生成的中间文件称为\underline{\hspace{4cm}}，其扩展名是\underline{\hspace{3cm}}。

	\item 面向对象的三大核心特征分别是：封装、\underline{\hspace{4cm}} 和 \underline{\hspace{4cm}}。

	\item 在Java的异常处理体系中，所有的异常类都继承自 \texttt{Throwable} 类，它主要有两个直接子类：\underline{\hspace{4cm}}（用于系统严重错误）和 \underline{\hspace{4cm}}（用于程序可以处理的异常）。

	\item Java中线程的生命周期包含了新建（New）、就绪（Runnable）、\underline{\hspace{3cm}}（Running）、\underline{\hspace{3cm}}（Blocked）以及死亡（Dead）五种基本状态。

	\item 在Java集合框架中，\underline{\hspace{4cm}} 接口的实现类允许存在重复的元素并且是有序的；而 \underline{\hspace{4cm}} 接口的实现类不允许存在重复的元素。
\end{enumerate}


\section*{四、阅读程序，写出运行结果（本大题共 5 小题，每小题 6 分，共 30 分）}

\begin{enumerate}[label=\arabic*., itemsep=1em]
	\item \textbf{写出下列程序的输出结果：}
	      \begin{lstlisting}[style=JavaExam]
public class Test1 {
    public static void main(String[] args) {
        int sum = 0;
        for (int i = 1; i <= 5; i++) {
            if (i == 3) continue;
            sum += i;
        }
        System.out.println("sum = " + sum);
    }
}
\end{lstlisting}
	      \textbf{输出结果：}\underline{\hspace{8cm}}

	\item \textbf{写出下列程序的输出结果：}
	      \begin{lstlisting}[style=JavaExam]
public class Test2 {
    public static void change(int[] arr) {
        arr[0] = 99;
    }
    public static void main(String[] args) {
        int[] nums = {10, 20, 30};
        change(nums);
        System.out.println(nums[0] + ", " + nums[1]);
    }
}
\end{lstlisting}
	      \textbf{输出结果：}\underline{\hspace{8cm}}

	\item \textbf{写出下列程序的输出结果：}
	      \begin{lstlisting}[style=JavaExam]
class Father {
    int x = 10;
    public void show() {
        System.out.println("Father Show");
    }
}
class Son extends Father {
    int x = 20;
    public void show() {
        System.out.println("Son Show");
    }
}
public class Test3 {
    public static void main(String[] args) {
        Father f = new Son();
        System.out.println(f.x);
        f.show();
    }
}
\end{lstlisting}
	      \textbf{输出结果：}\underline{\hspace{8cm}}

	\item \textbf{写出下列程序的输出结果：}
	      \begin{lstlisting}[style=JavaExam]
public class Test4 {
    public static void main(String[] args) {
        String s1 = new String("Java");
        String s2 = new String("Java");
        System.out.println(s1 == s2);
        System.out.println(s1.equals(s2));
    }
}
\end{lstlisting}
	      \textbf{输出结果：}\underline{\hspace{8cm}}

	\item \textbf{写出下列程序的输出结果：}
	      \begin{lstlisting}[style=JavaExam]
public class Test5 {
    public static int testException() {
        int res = 1;
        try {
            res = 10 / 0;
            return res;
        } catch (ArithmeticException e) {
            res = 2;
            return res;
        } finally {
            res = 3;
            System.out.print("Finally ");
        }
    }
    public static void main(String[] args) {
        System.out.println(testException());
    }
}
\end{lstlisting}
	      \textbf{输出结果：}\underline{\hspace{8cm}}
\end{enumerate}

\newpage
\section*{五、编程题（本大题共 2 小题，每小题 20 分，共 40 分）}

\textbf{1. 面向对象基础设计（20分）}\\
请编写一个表示银行账户的类 \texttt{BankAccount}，具体要求如下：
\begin{itemize}
	\item (1) 包含三个私有属性：账号 (\texttt{String accountId})、户名 (\texttt{String ownerName})、余额 (\texttt{double balance})。
	\item (2) 提供一个带有三个参数的构造方法，用于初始化这三个属性。
	\item (3) 提供一个存款方法 \texttt{public void deposit(double amount)}，调用时增加账户余额，并打印“存款成功”。
	\item (4) 提供一个取款方法 \texttt{public void withdraw(double amount)}，如果余额充足则扣除对应金额并打印“取款成功”，如果余额不足则不扣款并打印“余额不足”。
	\item (5) 提供一个显示信息的方法 \texttt{public void showInfo()}，输出账号、户名和当前余额。
\end{itemize}
\vspace{1em}
\textbf{【答题区】}
\vspace{12cm}

\textbf{2. 接口与多态的综合应用（20分）}\\
请按照以下要求编写程序：
\begin{itemize}
	\item (1) 定义一个接口 \texttt{Shape}，其中包含一个抽象方法 \texttt{double getArea()} 用于计算面积。
	\item (2) 编写一个实现 \texttt{Shape} 接口的类 \texttt{Rectangle}（矩形），包含私有属性 \texttt{width}（宽）和 \texttt{height}（高），提供构造方法，并重写 \texttt{getArea()} 方法。
	\item (3) 编写一个实现 \texttt{Shape} 接口的类 \texttt{Circle}（圆形），包含私有属性 \texttt{radius}（半径），提供构造方法，并重写 \texttt{getArea()} 方法（圆周率可使用 \texttt{Math.PI}）。
	\item (4) 编写一个测试类 \texttt{TestShape}，在 \texttt{main} 方法中创建一个大小为 2 的 \texttt{Shape} 数组。数组第一个元素存放一个宽为 4，高为 5 的矩形对象；第二个元素存放一个半径为 3 的圆形对象。使用 \texttt{for} 循环遍历数组，计算并输出这两个图形的面积之和。
\end{itemize}
\vspace{1em}
\textbf{【答题区】}

\newpage
% ---------------------------------------------------------
% 教师阅卷参考答案区（不需要打印给学生时可删除此页以下内容）
% ---------------------------------------------------------
\begin{center}
	{\huge\bfseries 教师阅卷参考答案及评分标准}
\end{center}
\vspace{1em}

\textbf{一、单项选择题（每小题2分，共40分）}
\begin{tabularx}{\textwidth}{XXXXX}
	1. B  & 2. C  & 3. B  & 4. B  & 5. B  \\
	6. C  & 7. C  & 8. D  & 9. B  & 10. D \\
	11. B & 12. B & 13. C & 14. C & 15. D \\
	16. B & 17. B & 18. C & 19. C & 20. B \\
\end{tabularx}

\vspace{1em}
\textbf{二、判断题（每小题2分，共20分）}
\begin{tabularx}{\textwidth}{XXXXX}
	1. $\checkmark$ & 2. $\times$ (未重写时equals也是比地址) & 3. $\checkmark$ & 4. $\times$     & 5. $\checkmark$         \\
	6. $\checkmark$ & 7. $\times$ (Exception应放在最后)  & 8. $\checkmark$ & 9. $\checkmark$ & 10. $\times$ (局部变量无默认值) \\
\end{tabularx}

\vspace{1em}
\textbf{三、填空题（每空2分，共20分）}
\begin{enumerate}[label=\arabic*., itemsep=0.5em]
	\item 字节码（或字节码文件） \quad \texttt{.class}
	\item 继承 \quad 多态 （顺序可换）
	\item \texttt{Error} \quad \texttt{Exception}
	\item 运行（状态） \quad 阻塞（状态）
	\item \texttt{List} \quad \texttt{Set}
\end{enumerate}

\vspace{1em}
\textbf{四、阅读程序写结果（每小题6分，共30分）}
\begin{enumerate}[label=\arabic*., itemsep=0.5em]
	\item \texttt{sum = 12}
	\item \texttt{99, 20}
	\item \texttt{10}\\ \texttt{Son Show} \quad (多态机制，属性看左边，方法看右边，错一行扣3分)
	\item \texttt{false}\\ \texttt{true} \quad (错一行扣3分)
	\item \texttt{Finally 2} \quad (先执行finally输出，最后返回catch中记录的2)
\end{enumerate}

\vspace{1em}
\textbf{五、编程题（每小题20分，共40分）}
\\
\textbf{第1题评分标准参考：}
定义属性且私有化（4分）；构造方法正确（4分）；deposit方法正确（4分）；withdraw方法逻辑正确（4分）；showInfo方法正确（4分）。
\begin{lstlisting}[style=JavaExam]
class BankAccount {
    private String accountId;
    private String ownerName;
    private double balance;

    public BankAccount(String accountId, String ownerName, double balance) {
        this.accountId = accountId;
        this.ownerName = ownerName;
        this.balance = balance;
    }

    public void deposit(double amount) {
        balance += amount;
        System.out.println("存款成功");
    }

    public void withdraw(double amount) {
        if (balance >= amount) {
            balance -= amount;
            System.out.println("取款成功");
        } else {
            System.out.println("余额不足");
        }
    }

    public void showInfo() {
        System.out.println("账号: " + accountId + ", 户名: " + ownerName + ", 余额: " + balance);
    }
}
\end{lstlisting}

\vspace{1em}
\textbf{第2题评分标准参考：}
Shape接口定义正确（3分）；Rectangle类完整正确（5分）；Circle类完整正确（5分）；测试类数组定义正确（3分）；循环及总面积计算输出正确（4分）。
\begin{lstlisting}[style=JavaExam]
interface Shape {
    double getArea();
}

class Rectangle implements Shape {
    private double width;
    private double height;
    public Rectangle(double width, double height) {
        this.width = width;
        this.height = height;
    }
    public double getArea() {
        return width * height;
    }
}

class Circle implements Shape {
    private double radius;
    public Circle(double radius) {
        this.radius = radius;
    }
    public double getArea() {
        return Math.PI * radius * radius;
    }
}

public class TestShape {
    public static void main(String[] args) {
        Shape[] shapes = new Shape[2];
        shapes[0] = new Rectangle(4, 5);
        shapes[1] = new Circle(3);
        
        double totalArea = 0;
        for (int i = 0; i < shapes.length; i++) {
            totalArea += shapes[i].getArea();
        }
        System.out.println("总面积为：" + totalArea);
    }
}
\end{lstlisting}

\end{document}